\section{Guarded extensionality and precision algebra}
\label{sec:extensionality}

\subsection{Observation equivalence rather than quotient types}

The object language of iPRA is first order over natural numbers.  It need not
contain quotient types for real numbers, probability measures, projective
points, or functions.  Instead, a mathematical object is handled through a
fixed numerical presentation and an observational equivalence relation.

For a functional presentation with evaluator \(a(p,k)\), define
\[
 p\equiv_X q
 \quad:\Longleftrightarrow\quad
 \forall k\;\Close_X(p,q,k),
\]
where \(\Close_X\) is a primitive-recursive finite inequality.  For a
relational presentation, the natural notion is Hausdorff-style finite
closeness of viable answer clouds, or equivalently agreement of all stable
finite tests.  In either case, the global equivalence is not inserted into an
induction formula.  At a fixed output challenge, only finitely many instances
of \(\Close_X\) are required.

\begin{principle}[Extensional proof discipline]
To prove that a numerical construction \(F\) is an operation on represented
objects, first choose a primitive-recursive precision transformer
\(\mu_F\).  Prove the fixed-precision implication
\[
 \bigwedge_{q\in Q_F(p,k)}\Close_X(p,p',q)
 \longrightarrow
 \Close_Y(F(p),F(p'),k),
\]
where \(Q_F(p,k)\) is the finite query table generated by \(\mu_F\) and the
bounded trace.  Only afterward introduce the universal quantifier over
\(k\).
\end{principle}

This makes extensionality a theorem with a visible numerical budget rather
than a hidden appeal to continuity.

\subsection{Composition of precision contracts}

A certified operation \(F:X\to Y\) has three numerical interfaces:

\begin{enumerate}
\item an output algorithm or output relation;
\item a demand function \(\mu_F(p,k)\), possibly returning a finite table of
      input challenges; and
\item an error transformer \(E_F\) proving that input errors within the
      demanded budget yield the requested output error.
\end{enumerate}

If \(G:Y\to Z\), the composite demand is obtained by first evaluating the
finite metadata needed for \(G\), then substituting every challenge requested
by \(G\) into \(F\)'s demand.  For scalar precision transformers this is the
familiar formula
\[
 \mu_{G\circ F}(p,k)=\mu_F(p,\mu_G(F_0(p),k)).
\]
For adaptive query tables, concatenate the proof-producing traces.  Error
budgets compose by ordinary rational inequalities.

\begin{theorem}[Precision-contract composition]
\label{thm:precision-composition}
Certified extensional transducers are closed under sequential composition,
finite products, and bounded case distinction.  The composite precision map,
trace bound, guard, and proof term are primitive-recursively computable from
the component data.
\end{theorem}

\begin{proof}
Sequential composition is the phase construction of the trace compiler.  A
finite product pairs states, unions the finite query tables, and takes the
maximum demanded precision.  A bounded case split computes every guard as a
Boolean and uses a phase tag to select the corresponding invariant.  In each
case the final extensionality implication is obtained by cut from the
component implications.  Construction of the resulting \(\Prf\) term is
structural recursion on the finite transducer expression.
\end{proof}

\subsection{Totalization and guards}

Primitive-recursive function symbols are total, whereas familiar
mathematical operations are often partial.  The correct solution is not to
pretend that the denominator, pivot, chart, logarithm argument, or derivative
is always valid.  Define an arbitrary totalized numerical operation and state
its correctness under a Boolean guard.

For example, a coded rational inverse may return zero when the numerator is
zero.  Its theorem has the form
\[
 \mathsf{Apart}(p,d)=1\to
 \mathsf{InvContract}(p,d,k,\mathsf{invApprox}(p,d,k))=1.
\]
The output outside the guard has no mathematical interpretation.  Because the
operation itself is total, it is a legitimate \(\PR\) symbol; because the
soundness theorem exposes the guard, no partiality is concealed.

\begin{principle}[No unguarded partial operation]
Every use of division, normalization, pivoting, branch inversion, logarithm,
argument, local chart, or reciprocal must consume a finite positive margin.
The margin may be an input certificate or may be produced by an earlier
verified module.  A mere negated equality is not silently converted to a
positive rational lower bound.
\end{principle}

\subsection{Overlapping branches}

Exact branching on an approximate comparison is unstable.  A robust
alternative uses overlap.  Suppose two formulas \(F_0(x)\) and \(F_1(x)\) are
both valid on an overlap region
\([	au-\eta,\tau+\eta]\), while the first is valid to the left and the second
to the right.  Approximate \(x\) within \(\eta/2\):

\begin{itemize}
\item if the approximation is below \(\tau\), return branch zero;
\item otherwise return branch one.
\end{itemize}

If the comparison is wrong because \(x\) lies near \(\tau\), both branches are
valid.  The construction is total and continuous at the level of its
contract, even though the returned branch bit may fluctuate.

\begin{theorem}[Robust overlapping branch]
\label{thm:overlapping-branch}
Given finite proofs of left validity, right validity, overlap validity, and an
approximation error below half the overlap width, iPRA proves that the branch
selected by the rational comparison satisfies the common output contract.
\end{theorem}

\begin{proof}
There are two decidable rational cases.  In the selected-left case, either the
true value is left of the overlap endpoint, where left validity applies, or
the approximation error places it in the overlap.  The selected-right case is
symmetric.  The proof is finite rational arithmetic and can be Booleanized.
\end{proof}

\subsection{Exact branching under a gap}

If an input certificate gives
\[
 x\le \tau-g\quad\text{or}\quad x\ge \tau+g
\]
for a positive rational \(g\), an approximation within \(g/3\) determines the
correct side exactly.  Thus exact sign, orientation, chart, pivot, and root
labels are available under quantitative separation.

\begin{corollary}[Exact branch from a positive gap]
\label{cor:exact-branch-gap}
A supplied disjunctive gap certificate and a sufficiently accurate rational
observation yield a primitive-recursive exact branch selector whose
correctness is provable in iPRA.
\end{corollary}

The disjunction is important.  From \(x\ne0\) alone, intuitionistic logic does
not provide either a sign or a rational lower bound for \(|x|\).

\subsection{Congruence for relational outputs}

A relational operation is extensional when close input clouds produce close
output clouds in both directions: every valid output from one input has a
compatible nearby output for the other, and conversely.  This is a finite
back-and-forth condition at each challenge.  It composes by witness
amalgamation.

\begin{theorem}[Relational congruence]
\label{thm:relational-congruence}
Let \(F:X\rightsquigarrow Y\) be a certified relational transducer.  Suppose
input parameters \(p,p'\) satisfy the finite closeness table requested for an
output challenge \(k\).  Then every accepted output answer for \(p\) has a
\(Y\)-compatible accepted output answer for \(p'\), and conversely, provided
the transducer's local back-and-forth obligations are derivable.
\end{theorem}

\begin{proof}
Run the two finite traces in product.  Add the input closeness facts to the
Boolean invariant and pair each nondeterministic answer with a finite set of
compatible answers on the other side.  Backward viability removes dead ends.
The final product invariant gives the two output witnesses.  All searches are
bounded, so the argument is an instance of finite witness amalgamation and
embedded trace reflection.
\end{proof}

This theorem is the appropriate substitute for quotient congruence when
objects have non-canonical representatives.  A deterministic operation is the
special case in which each output cloud is a singleton.

\subsection{Normalization is a theorem, not an identity assumption}

Codes for rationals, polynomials, projective points, matrices, and finite sets
often admit multiple representations.  The framework uses one of two safe
patterns:

\begin{enumerate}
\item choose a primitive-recursive canonical normalization and prove that it
      preserves observations; or
\item retain non-canonical codes and prove every operation respects the
      observational relation.
\end{enumerate}

The first pattern is convenient for signed integers and positive-denominator
rationals.  The second is often better for projective or clustered data, where
canonical choice would require unstable comparisons.  In neither pattern is
code equality confused with mathematical equality.
